% =========================================================================
% ALGORITMO NETWORK SIMPLEX E ÁRVORES GERADORAS
% =========================================================================
\chapter{A Abordagem Primal-Dual em Árvore: O Algoritmo Network Simplex}\label{sec:network_simplex}

\section{Árvore Geradora Básica e Partição de Arcos}

Uma solução básica viável do MCF é caracterizada por uma partição dos arcos do digrafo em três conjuntos disjuntos (veja~\cite{ahuja1993network}):
\begin{enumerate}
    \item \textbf{Arcos básicos} ($T$): formam uma árvore geradora do digrafo. O fluxo nestes arcos pode assumir qualquer valor entre $0$ e $c(u,v)$.
    \item \textbf{Arcos não básicos no limite inferior} ($L$): arcos fora da árvore com fluxo fixado em $f(u,v) = 0$.
    \item \textbf{Arcos não básicos no limite superior} ($U$): arcos fora da árvore com fluxo fixado em $f(u,v) = c(u,v)$ (saturados).
\end{enumerate}

Dada uma árvore geradora $T$, o fluxo nos arcos básicos é unicamente determinado pelas demandas dos vértices e pelos fluxos fixados nos arcos de $L$ e $U$ (através da conservação de fluxo). Da mesma forma, os potenciais nodais $\pi$ são determinados univocamente (a menos de uma constante aditiva) pela condição de que $\bar{w}_\pi(u,v) = 0$ para todo arco $(u,v) \in T$, ou seja:
\begin{equation}
    \pi(u) - \pi(v) = w(u,v), \quad \forall (u,v) \in T
\end{equation}

Fixando $\pi(r) = 0$ para um vértice raiz $r$ arbitrário, os potenciais de todos os demais vértices são propagados pela árvore em tempo $\mathcal{O}(|V|)$.

\begin{figure}[!ht]
\centering
\begin{tikzpicture}[thick]

    \node[source node] (v1) at (0, 0) {$v_1$};
    \node[flow node] (v2) at (2.5, 1.5) {$v_2$};
    \node[flow node] (v3) at (2.5, -1.5) {$v_3$};
    \node[flow node] (v4) at (5, 1.5) {$v_4$};
    \node[sink node] (v5) at (5, -1.5) {$v_5$};

    \draw[tree edge] (v1) -- node[above left,font=\footnotesize] {$T$} (v2);
    \draw[tree edge] (v1) -- node[below left,font=\footnotesize] {$T$} (v3);
    \draw[tree edge] (v2) -- node[above,font=\footnotesize] {$T$} (v4);
    \draw[tree edge] (v3) -- node[below,font=\footnotesize] {$T$} (v5);

    \draw[nontree edge] (v2) -- node[edge lbl right] {$L$} (v3);
    \draw[nontree edge] (v4) -- node[edge lbl right] {$U$} (v5);
\end{tikzpicture}
\caption{Partição dos arcos no Network Simplex. Arcos azuis sólidos formam a árvore geradora básica $T$ (4 arcos para 5 vértices). Arcos cinza tracejados são não-básicos: $L$ (fluxo $= 0$) e $U$ (fluxo $= c$, saturado).}
\label{fig:arvore_base}
\end{figure}


\section{Estrutura de Dados da Árvore e Inicialização com Raiz Artificial}

Para viabilizar a execução eficiente das travessias em tempo $\mathcal{O}(|V|)$, a implementação do \textit{Network Simplex} (Código~\ref{lst:network_simplex}, Apêndice~\ref{ap:network_simplex}) mantém a árvore geradora representada por três vetores principais:
\begin{itemize}
    \item \lstinline{parent[v]}: armazena o vértice predecessor direto de $v$ na árvore direcionada.
    \item \lstinline{depth[v]}: registra a profundidade do vértice em relação à raiz na árvore geradora, facilitando a identificação do ancestral comum mais próximo (\textit{LCA - Lowest Common Ancestor}) no ciclo fundamental.
    \item \lstinline{thread[v]}: mantém uma lista encadeada da travessia em pré-ordem da árvore, permitindo iterar sobre toda a subárvore enraizada em $v$ em tempo proporcional ao número de seus descendentes durante a atualização dos potenciais nodais.
\end{itemize}

\textbf{Construção da Base Inicial:} Para garantir que uma árvore geradora viável inicial esteja sempre disponível sem necessidade de computação preliminar complexa, introduz-se um vértice raiz artificial $r = |V|$. Criam-se arcos artificiais conectando a raiz a cada vértice original com capacidade infinita e custo unitário suficientemente grande $M$ (\textit{Big-M Method}). Estes arcos formam a árvore inicial $T$. Ao longo das iterações do Simplex, os arcos originais da rede entram na base substituindo progressivamente os arcos artificiais, terminando em uma solução composta exclusivamente pelos arcos do digrafo original.


\section{Total Unimodularidade (TUM) e Soluções Inteiras}\label{sec:tum_secao}

A razão algébrica pela qual o \textit{Network Simplex} sempre produz soluções inteiras (quando capacidades e demandas são inteiras) reside em uma propriedade fundamental da matriz de restrições do problema de fluxo: a \textbf{Total Unimodularidade} (TUM).

\begin{quote}
\textbf{Definição.} Uma matriz $M \in \mathbb{Z}^{m \times n}$ é \textit{Totalmente Unimodular} (TUM) se o determinante de toda submatriz quadrada de $M$ pertence ao conjunto $\{-1, 0, +1\}$.
\end{quote}

A matriz de restrições do problema de fluxo em redes é a \textbf{matriz de incidência} $A \in \{-1, 0, +1\}^{|V| \times |A|}$ do digrafo $D = (V, A)$, definida como:
\begin{equation}
    A_{v,e} = \begin{cases}
        +1, & \text{se } v \text{ é a origem do arco } e \\
        -1, & \text{se } v \text{ é o destino do arco } e \\
        \phantom{+}0, & \text{caso contrário}
    \end{cases}
\end{equation}

\begin{teorema}[Total Unimodularidade da Matriz de Incidência --- Poincaré, 1901][teo:tum]
A matriz de incidência nó-arco $A$ de qualquer digrafo direcionado é Totalmente Unimodular (TUM), isto é, o determinante de qualquer submatriz quadrada de $A$ pertence ao conjunto $\{-1, 0, +1\}$.
\end{teorema}

\intuicao{Integralidade Poliedral} Pela Regra de Cramer, as componentes de uma solução básica $B x_B = b$ satisfazem $x_j = \det(B_j) / \det(B)$. Sendo $A$ totalmente unimodular e $B$ uma base viável, tem-se $\det(B) \in \{-1, +1\}$. A divisão por $\pm 1$ preserva a integralidade de qualquer vetor de demandas inteiras $b \in \mathbb{Z}^{|V|}$, assegurando que todos os vértices do poliedro de fluxos sejam estritamente inteiros.

\begin{demonstracao}
Procedemos por indução na ordem $k$ da submatriz quadrada arbitrária $M \in \mathbb{R}^{k \times k}$ formada a partir de $A$:

\medskip
\noindent \textit{Base de indução ($k = 1$):} \\
Toda submatriz de ordem $1$ consiste em uma única entrada escalar $A_{v,e} \in \{-1, 0, +1\}$, logo $\det(M) \in \{-1, 0, +1\}$.

\medskip
\noindent \textit{Passo indutivo:} \\
Assuma como hipótese que qualquer submatriz quadrada de ordem $k-1$ possua determinante em $\{-1, 0, +1\}$. Seja $M$ uma submatriz quadrada de ordem $k > 1$. Como cada coluna de $A$ possui no máximo um elemento $+1$ e um $-1$, cada coluna de $M$ contém no máximo dois elementos não nulos:

\begin{enumerate}
    \item \textbf{Existe coluna identicamente nula em $M$:} $\det(M) = 0 \in \{-1, 0, +1\}$.
    
    \item \textbf{Existe coluna com exatamente um elemento não nulo:} \\
    Seja a $j$-ésima coluna contendo apenas $M_{i,j} \in \{-1, +1\}$ na linha $i$. Pela expansão de Laplace por cofatores:
    $$\det(M) = M_{i,j} \cdot (-1)^{i+j} \cdot \det(M_{i,j})$$
    onde $M_{i,j}$ é de ordem $k-1$. Por hipótese de indução, $\det(M_{i,j}) \in \{-1, 0, +1\}$, logo $\det(M) \in \{-1, 0, +1\}$.
    
    \item \textbf{Todas as colunas possuem exatamente dois elementos não nulos:} \\
    Cada coluna contém exatamente um $+1$ e um $-1$. Somando todas as linhas de $M$, a soma em cada coluna anula-se:
    $$\sum_{i=1}^k M_{i,j} = (+1) + (-1) = 0, \quad \forall\, j \in \{1, \dots, k\}$$
    Logo, $\mathbf{1}^T M = \mathbf{0}^T$, as linhas de $M$ são linearmente dependentes e $\det(M) = 0 \in \{-1, 0, +1\}$.
\end{enumerate}
Por indução, a matriz de incidência $A$ é Totalmente Unimodular.
\end{demonstracao}

\implicacoes{para a Estrutura Poliédrica e o Network Simplex} A total unimodularidade da matriz de incidência assegura que problemas de fluxo linear com parâmetros inteiros possuam soluções ótimas inteiras sem restrições explícitas de integralidade ($\mathbb{Z}$), eliminando a necessidade de métodos de corte ou branch-and-bound.


\section{Potenciais Nodais e Custos Reduzidos na Árvore}

Uma vez determinada a árvore $T$ e fixado $\pi(r) = 0$, os potenciais são calculados por travessia na árvore, aplicando recursivamente $\pi(v) = \pi(u) - w(u,v)$ para cada arco $(u,v) \in T$.

Com os potenciais calculados, uma solução básica é ótima se e somente se:
\begin{itemize}
    \item $\bar{w}_\pi(u,v) \ge 0$, para todo $(u,v) \in L$.
    \item $\bar{w}_\pi(u,v) \le 0$, para todo $(u,v) \in U$.
\end{itemize}

Se algum arco viola essas condições, ele é candidato a entrar na base, desencadeando uma iteração do Simplex.


\section{Iteração do Network Simplex: Pricing, Ciclo Fundamental e Pivoteamento}

Cada iteração do \textit{Network Simplex} compreende quatro etapas~\cite{ahuja1993network, dantzig1951application}:

\textbf{Passo 1 --- Pricing (Seleção do arco entrante):}
Identifica-se um arco não básico que viola a otimalidade: $(u,v) \in L$ com $\bar{w}_\pi(u,v) < 0$ ou $(u,v) \in U$ com $\bar{w}_\pi(u,v) > 0$. A implementação utiliza busca por blocos (\textit{Block Search Pricing}) para acelerar a busca. Se não houver violações, a solução é ótima.

\textbf{Passo 2 --- Ciclo Fundamental:}
A adição do arco entrante $(u,v)$ à árvore $T$ forma um ciclo fundamental $C$, identificado por busca até o ancestral comum mais próximo (\textit{LCA}).

\textbf{Passo 3 --- Arco limitante (\textit{leaving arc}):}
O gargalo $\delta$ é a menor folga de capacidade ao longo de $C$. O arco que atinge seu limite deixa a base. Para evitar ciclos infinitos sob degenerescência ($\delta = 0$), utiliza-se a regra de base fortemente viável de Cunningham~\cite{cunningham1976network}.

\textbf{Caso Especial (Arco Auto-Limitante):} Quando $\delta = c(e_{in})$, o arco entrante atinge seu limite antes de qualquer arco básico. A implementação alterna o estado de $e_{in}$ entre $L$ e $U$ e satura o fluxo, sem alterar a árvore $T$.

\textbf{Passo 4 --- Pivoteamento e Atualização da Árvore em $\mathcal{O}(|V|)$:}
O fluxo é atualizado por $\delta$ ao longo de $C$. O arco $e_{out}$ sai da base e $e_{in}$ entra em $T$. A remoção de $e_{out}$ particiona a árvore em $T_r$ (contendo a raiz $r$) e $T_s$. Os potenciais em $T_s$ são atualizados por:
$$\pi(x) \leftarrow \pi(x) \pm \bar{w}_\pi(e_{in}), \quad \forall x \in V(T_s)$$
A lista encadeada da travessia em pré-ordem (\lstinline{thread}) permite atualizar os nós de $T_s$ em $\mathcal{O}(|V|)$. Os vetores \lstinline{parent} e \lstinline{depth} são ajustados ao longo do caminho até a raiz de $T_s$.

\begin{figure}[!ht]
\centering
\begin{tikzpicture}[thick]

\begin{scope}[xshift=0cm]
  \tikzsubcaption{(2.0,4.2)}{(a) Antes do pivot}
  \node[source node] (a1) at (0, 1.5) {$1$};
  \node[flow node] (a2) at (2, 3.0) {$2$};
  \node[flow node] (a3) at (2, 0.0) {$3$};
  \node[sink node] (a4) at (4, 1.5) {$4$};

  \draw[tree edge] (a1) -- (a2);
  \draw[tree edge] (a2) -- (a4);
  \draw[tree edge] (a1) -- (a3);
  \draw[pivot enter] (a3) -- node[below=3.8pt,sloped,font=\footnotesize,text=red!80!black] {entra} (a4);
\end{scope}

\node at (5.5,1.5) {\Large $\Rightarrow$};

\begin{scope}[xshift=7cm]
  \tikzsubcaption{(2.0,4.2)}{(b) Depois do pivot}
  \node[source node] (b1) at (0, 1.5) {$1$};
  \node[flow node] (b2) at (2, 3.0) {$2$};
  \node[flow node] (b3) at (2, 0.0) {$3$};
  \node[sink node] (b4) at (4, 1.5) {$4$};

  \draw[tree edge] (b1) -- (b2);
  \draw[pivot leave] (b2) -- node[above=3.8pt,sloped,font=\footnotesize,text=gray!80!black] {sai} (b4);
  \draw[tree edge] (b1) -- (b3);
  \draw[tree edge] (b3) -- (b4);
\end{scope}
\end{tikzpicture}
\caption{Iteração do Network Simplex: (a) o arco $(3,4)$ entra na base, fechando um ciclo fundamental. (b) Após o incremento $\delta$, o arco $(2,4)$ sai da base e $(3,4)$ passa a integrar a árvore geradora.}
\label{fig:simplex_pivoteamento}
\end{figure}

\paragraph{Complexidade e Desempenho Prático.}
Embora o Network Simplex não apresente garantia de complexidade polinomial no pior caso sob regras de pivoteamento usuais, cada pivoteamento opera em $\mathcal{O}(|V|)$ e o algoritmo executa um número reduzido de iterações em instâncias práticas da literatura~\cite{ahuja1993network}. A implementação em C++ está detalhada no Código~\ref{lst:network_simplex} (Apêndice~\ref{ap:network_simplex}).


